\documentclass[11pt,oneside]{report}

\usepackage[a4paper,margin=25mm,headheight=15pt]{geometry}
\usepackage{fontspec}
\usepackage{unicode-math}
\setmainfont{FreeSerif}
\setsansfont{FreeSans}
\setmonofont{FreeMono}[Scale=0.86]
\setmathfont{Asana Math}
\usepackage{amsmath,mathtools}
\usepackage{booktabs,longtable,tabularx,array}
\usepackage{xcolor}
\usepackage{microtype}
\usepackage{fancyhdr}
\usepackage{titlesec}
\usepackage[most]{tcolorbox}
\usepackage[
  colorlinks=true,
  linkcolor=blue!55!black,
  urlcolor=blue!60!black,
  citecolor=green!35!black
]{hyperref}
\usepackage{bookmark}
\usepackage[
  hybrid,
  pipeTables,
  tableCaptions,
  fencedCode,
  smartEllipses
]{markdown}

\definecolor{ink}{HTML}{18212B}
\definecolor{accent}{HTML}{315D8A}
\definecolor{softblue}{HTML}{EEF5FB}
\definecolor{softgold}{HTML}{FFF7E8}
\definecolor{deepgreen}{HTML}{245A45}
\color{ink}
\setlength{\parindent}{0pt}
\setlength{\parskip}{0.55em}
\setcounter{tocdepth}{2}
\setcounter{secnumdepth}{2}
\emergencystretch=3em
\raggedbottom

\titleformat{\chapter}[display]
  {\normalfont\huge\bfseries\color{accent}}
  {\filleft\Large\chaptername\ \thechapter}
  {0.8ex}
  {\titlerule\vspace{1ex}\filright}
\titleformat{\section}
  {\normalfont\Large\bfseries\color{accent}}
  {\thesection}{0.6em}{}
\titleformat{\subsection}
  {\normalfont\large\bfseries\color{deepgreen}}
  {\thesubsection}{0.6em}{}

\pagestyle{fancy}
\fancyhf{}
\fancyhead[L]{\small Mean-Field Peeling}
\fancyhead[R]{\small Self-Contained Research Report}
\fancyfoot[C]{\thepage}

% Preserve Markdown equation labels that occur inside $$...$$.
\renewcommand{\tag}[1]{\qquad\text{\normalfont (#1)}}

\newtcolorbox{statusbox}{
  colback=softblue,colframe=accent,boxrule=0.7pt,arc=2mm,
  left=2mm,right=2mm,top=1.5mm,bottom=1.5mm
}
\newtcolorbox{warningbox}{
  colback=softgold,colframe=orange!65!black,boxrule=0.7pt,arc=2mm,
  left=2mm,right=2mm,top=1.5mm,bottom=1.5mm
}
\newcommand{\odotprod}{\mathbin{\odot}}

\begin{document}

\begin{titlepage}
\centering
\vspace*{2.2cm}
{\Huge\bfseries\color{accent} Mean-Field Peeling\par}
\vspace{0.45cm}
{\LARGE A Self-Contained Research Report\par}
\vspace{0.65cm}
{\Large Explicit Gaussian Normal Forms, Backward Kernels,\\
and Local Feature-Learning Jets\par}
\vfill
\begin{tcolorbox}[width=0.86\textwidth,colback=softblue,colframe=accent]
\centering
\textbf{Purpose.}
This report consolidates the original proposal, the maintained mathematical
interpretation, the audited backward-kernel execution, the one-step and
two-step $\mu$P calculations, and the current theorem and literature
positioning.  It is designed for uninterrupted offline reading.
\end{tcolorbox}
\vfill
{\large Amir Joudaki's mean-field peeling research program\par}
\vspace{0.25cm}
{\large Consolidated and audited 3 August 2026\par}
\vspace{1.2cm}
\end{titlepage}

\pagenumbering{roman}
\tableofcontents
\clearpage

\chapter*{Preface: what is established and what remains open}
\addcontentsline{toc}{chapter}{Preface: status of the theory}

Mean-field derivative peeling is a proposed layerwise elimination calculus
for fixed-order scalar derivative observables of wide neural networks.  Its
desired output is a \emph{Gaussian normal form}: a finite directed acyclic
graph whose leaves are fully specified Gaussian expectations and whose
remaining operations are deterministic algebra.

\begin{statusbox}
\textbf{Current status.}
The conditional Gaussian identities, finite-width chain-rule expansions,
backward-kernel expectation calculation, $\mu$P scaling, one-step tangent
factorization, exact two-Euler-step Taylor factors, and deep-linear audits are
explicitly derived.  The general normal-form theorem, a complete concentration
theorem, depth-uniform finite-state closure, and a fully enumerated nonlinear
two-step channel registry remain proof obligations.
\end{statusbox}

The report therefore separates exact finite-width identities, mean-field
results under explicit assumptions, formal recursions, audited restricted
calculations, target theorems, and open implications.

\clearpage
\pagenumbering{arabic}

\part{A Guided Synthesis}

\chapter{The central problem}

Fix a batch of $B$ input--label pairs and a fully connected network of fixed
depth $L$ and width $n$:
\[
h^0(a)=x_a,\qquad
z^\ell(a)=W^\ell h^{\ell-1}(a),\qquad
h^\ell(a)=\phi(z^\ell(a)).
\]
The objects of interest are scalar contractions built from outputs, losses,
hidden Grams, and a fixed number of parameter or training derivatives.  A
representative family is
\[
[\eta^r]\,
\mathbb E\!\left[
\frac1n\langle h^{\ell,(s)}_a(\eta),
h^{\ell,(s)}_b(\eta)\rangle
\right],
\]
where derivative order $r$, step count $s$, batch size, depth, and observable
syntax are fixed before $n\to\infty$.

The proposed calculation is:
\begin{enumerate}
\item expand every derivative into products of explicit weights,
preactivations, and scalar activation derivatives;
\item expose every neuron and batch index;
\item select the highest active layer;
\item condition on lower layers and integrate the current Gaussian matrix;
\item execute every Stein/Onsager correction;
\item enumerate index-equality patterns and count their global width powers;
\item replace surviving row averages by deterministic Gaussian moments only
after every index using them is exposed; and
\item pass the resulting boundary state to the next lower layer.
\end{enumerate}
After the last peel, no random weight or width index should remain.

\chapter{The exact local calculus}

\section{Conditional Gaussianity}

Global joint Gaussianity of all deep variables is false.  The reusable exact
fact is layerwise conditional Gaussianity.  Let
\[
w\sim N(0,I_n),\qquad
Z_a=\frac1{\sqrt n}\sum_jw_jh_j(a).
\]
Conditional on the lower-layer values $h_j(a)$, $(w,Z)$ is jointly Gaussian
and
\[
\operatorname{Cov}(w_p,Z_a\mid h)
=\frac{h_p(a)}{\sqrt n}.
\]

\section{One- and two-weight identities}

For smooth $F:\mathbb R^B\to\mathbb R$,
\[
\mathbb E[w_pF(Z)\mid h]
=
\frac1{\sqrt n}\sum_a h_p(a)
\mathbb E[\partial_aF(Z)\mid h],
\]
and
\[
\mathbb E[w_pw_qF(Z)\mid h]
=
\delta_{pq}\mathbb E[F(Z)\mid h]
+
\frac1n\sum_{a,b}h_p(a)h_q(b)
\mathbb E[\partial_a\partial_bF(Z)\mid h].
\]
The first branch is Wick pairing.  The second is the complete Stein/Onsager
correction.  It cannot be deleted before global width counting.

\section{Index ownership and mean-field replacement}

An index may be summed during group $m$ only when it appears in no lower
group.  A factor created by Stein differentiation belongs to the group of its
preactivation source and becomes boundary data for the next peel.

The exact conditional law initially uses empirical covariances such as
\[
G_{n,ab}^{\ell-1}
=\frac1n\sum_jh_j^{\ell-1}(a)h_j^{\ell-1}(b).
\]
One first performs conditional integration by parts, retains all lower-group
factors, exposes equality patterns, and only then replaces $G_n^{\ell-1}$ by
its deterministic limit.  Distinguished rows require leave-one-out or
leave-two-out covariances and uniform integrability.

\chapter{The theorem program}

\section{Conservative target}

An admissible observable must be defined syntactically from outputs, losses,
preactivations, activations, hidden Grams, fixed-order parameter derivatives,
finite products, sums, and complete tensor contractions, with every
normalization declared.

\begin{statusbox}
\textbf{Target Mean-Field Derivative Peeling Theorem.}
For every admissible scalar derivative observable $O_n$, a deterministic
transformation recursively eliminates the maximal active layer.  Each peel
emits explicit Gaussian-integral coefficients multiplying admissible
lower-group states.  Under smoothness, moment, leave-one-out,
conditional-CLT, concentration, and uniform-integrability hypotheses,
\[
\lim_{n\to\infty}\mathbb E[O_n]
=\operatorname{Eval}(\operatorname{Peel}(O_n)).
\]
Every covariance, integrand, activation derivative, equality pattern, and
coefficient in the output is mechanically specified.
\end{statusbox}

This is a target theorem, not a result proved by the current examples.

\section{Finite-state strengthening}

The strongest useful extension would prove that, at fixed batch size,
derivative order, step count, and observable schema, the boundary signatures
close on finitely many depth-independent state types.  After sharing common
states, the normal-form DAG would then have at most
\[
C_{B,r,O}L
\]
nodes.  The present nonlinear registry is not fully enumerated and its channel
count may depend on depth.  Thus $O(L)$ complexity remains an open target.

\section{Feature-learning jets}

For $s$ training steps and Taylor order $p$, the coefficients
\[
[\eta^k]\mathbb E[\mathcal L_a^{(s)}(\eta)],
\qquad
[\eta^k]\mathbb E[G_{ab}^{\ell,(s)}(\eta)],
\qquad k\le p,
\]
form an order-$p$ \emph{feature-learning jet}.  A peeling theorem would give
each admissible coefficient an explicit Gaussian normal form.

\begin{warningbox}
A fixed jet does not imply convergence of the Taylor series, control at a
fixed nonzero learning rate, or a complete positive-time training trajectory.
\end{warningbox}

\chapter{The audited μP laboratory}

\section{Network and kernels}

The worked model has three hidden layers:
\[
z_i^1(a)=\frac1{\sqrt{d_0}}\sum_jW^1_{ij}x_{aj},\quad
z_p^2(a)=\frac1{\sqrt n}\sum_iW^2_{pi}h_i^1(a),
\]
\[
z_r^3(a)=\frac1{\sqrt n}\sum_pW^3_{rp}h_p^2(a),\quad
f_a=\frac1n\sum_rv_rh_r^3(a),
\]
with iid standard-Gaussian raw parameters and
\[
\mathcal J=\frac1B\sum_c(f_c-y_c)^2,\qquad
\theta^+=\theta-n\eta\nabla_\theta\mathcal J.
\]
Let
\[
G^0_{ab}=\frac{x_a^\top x_b}{d_0},\qquad
G^\ell_{ab}=\mathbb E[\phi(Z^\ell_a)\phi(Z^\ell_b)],
\]
\[
D^\ell_{ab}=\mathbb E[\phi'(Z^\ell_a)\phi'(Z^\ell_b)],
\qquad Z^\ell\sim N(0,G^{\ell-1}).
\]

\section{Backward sensitivity}

Define
\[
\delta_i^\ell(a)=n\frac{\partial f_a}{\partial z_i^\ell(a)},\qquad
\Pi^\ell_{n,ab}=\frac1n\sum_i\delta_i^\ell(a)\delta_i^\ell(b).
\]
The top-down Wick--Stein calculation gives the expectation limits
\[
\Pi^3=D^3,\qquad
\Pi^2=D^2\odotprod D^3,\qquad
\Pi^1=D^1\odotprod D^2\odotprod D^3.
\]
Random-kernel convergence requires a separate four-path concentration peel.

\section{One training step}

The optimizer-metric tangent kernel is exactly
\[
K_{n,ab}
=G^3_{n,ab}
+G^2_{n,ab}\Pi^3_{n,ab}
+G^1_{n,ab}\Pi^2_{n,ab}
+G^0_{ab}\Pi^1_{n,ab}.
\]
Under the joint concentration assumptions,
\[
\begin{aligned}
K_{ab}
={}&G^3_{ab}+G^2_{ab}D^3_{ab}
+G^1_{ab}D^2_{ab}D^3_{ab}\\
&+G^0_{ab}D^1_{ab}D^2_{ab}D^3_{ab}.
\end{aligned}
\]
Consequently,
\[
\lim_{n\to\infty}[\eta]\,
\mathbb E[(f_a^+-y_a)^2]
=-\frac4B y_a\sum_cK_{ac}y_c.
\]
The annealed hidden-Gram linear coefficient vanishes,
\[
\lim_{n\to\infty}[\eta]\,
\mathbb E[G_{n,ab}^{\ell,+}]=0,
\]
while the first generic deterministic coefficient is quadratic:
\[
\lim_{n\to\infty}[\eta^2]\,
\mathbb E[G_{n,ab}^{\ell,+}]=C^\ell_{ab}.
\]
The detailed recursion contains a fresh off-diagonal Gaussian field; replacing
the empirical random operator by its mean before squaring gives a wrong
coefficient.

\section{Two Euler steps}

For $F(\theta)=-n\nabla\mathcal J(\theta)$,
\[
\theta^{(2)}
=\theta+2\eta F+\eta^2DF[F]+O(\eta^3).
\]
For a hidden-Gram observable $O$, put
\[
L=DO[F],\qquad
C=\frac12D^2O[F,F],\qquad
R=DO[DF[F]].
\]
Then the cumulative and second-update expansions are
\[
O(\theta^{(2)})
=O(\theta)+2\eta L+\eta^2(4C+R)+O(\eta^3),
\]
\[
O(\theta^{(2)})-O(\theta^{(1)})
=\eta L+\eta^2(3C+R)+O(\eta^3).
\]
The factors $4$ and $3$ are exact.  The general nonlinear reduction of $R$
is a formal multichannel schema pending a complete registry and convergence
theorem.

\section{Deep-linear audit}

For $\phi(z)=z$, with $\kappa=2/B$ and $\sigma=G^0y$,
\[
(C^1,C^2,C^3)_{ab}
=(1,5,14)\kappa^2\sigma_a\sigma_b,
\]
\[
(R^1,R^2,R^3)_{ab}
=(6,14,20)\kappa^2\sigma_a\sigma_b.
\]
Thus the cumulative two-step coefficients are $(10,34,76)$ and the
second-update coefficients are $(9,29,62)$, each multiplying
$\kappa^2\sigma_a\sigma_b$.

\chapter{Lessons, positioning, and reading route}

\section{Corrections forced by execution}

\begin{enumerate}
\item Fix raw versus effective parameter conventions before counting width.
\item Peel from the highest group downward.
\item Perform exact conditional Stein integration before covariance
replacement.
\item Retain Stein-created lower-group factors as boundary data.
\item Do not discard a zero-mean family before checking a second copy.
\item Separate expectation limits from concentration.
\item Do not differentiate a limiting zero as though it were identically zero.
\item At later steps, track differentiated states and every cross-covariance.
\item Treat network states as nonlinear functions of Gaussian base channels,
not as one jointly Gaussian state.
\item Use a safe fixed multiple $qB$ for Gaussian dimension; exactly $B$ is
not established generally.
\end{enumerate}

\section{Position relative to Tensor Programs}

Tensor Programs already gives constructive recursive limiting semantics for
broad finite neural computations, including $\mu$P feature-learning limits
and matrix-reuse Onsager corrections.  The defensible proposed contribution of
peeling is narrower: a source-to-source normal-form compiler for contracted
derivative observables that explicitly executes every derivative expansion,
Stein correction, Wick contraction, equality partition, and width count.

\section{How to use the detailed parts}

Part II preserves both the original notes and the maintained mathematical
program.  Part III is the audited calculation-level case study.  Part IV
states the sharp theorem target, the Tensor Programs contrast, required
nonclaims, and the decisive remaining proof obligations.  Repetition across
these parts is intentional: each preserves a distinct layer of the research
state.

% The embedded Markdown sources contain mathematics in headings. Starred
% headings preserve hierarchy without forcing the PDF bookmark machinery to
% serialize arbitrary TeX mathematics.
\markdownSetup{rendererPrototypes={
  headingOne={\section*{#1}},
  headingTwo={\subsection*{#1}},
  headingThree={\subsubsection*{#1}},
  headingFour={\paragraph*{#1}\leavevmode},
  headingFive={\subparagraph*{#1}\leavevmode}
}}
\input{studies/derivative_wick_calculus/build/markdown/markdown_math_defs.tex}

\part{The Maintained Peeling Program}
\chapter{Original notes, current interpretation, and proof program}
\markdownInput{studies/derivative_wick_calculus/build/markdown/MEAN_FIELD_PEELING_PROGRAM.pdf.md}

\part{Audited Worked Case Study}
\chapter{Backward kernels and one-step/two-step μP feature learning}
\markdownInput{studies/derivative_wick_calculus/build/markdown/MUP_TRAINING_PEELING_CASE_STUDY.pdf.md}

\part{Theorem and Literature Positioning}
\chapter{Tensor Programs contrast and the sharp theorem claim}
\markdownInput{studies/derivative_wick_calculus/build/markdown/TENSOR_PROGRAMS_CONTRAST_AND_THEOREM_FRAMING_REPORT.pdf.md}

\end{document}
